% Generated draft: requires mathematical and source review.
\documentclass[11pt]{article}
\usepackage[margin=25mm]{geometry}
\usepackage{fontspec}
\setmainfont{Latin Modern Roman}
\usepackage{amsmath,amssymb,mathtools,unicode-math}
\setmathfont{Latin Modern Math}
\usepackage{enumitem,needspace,xurl,hyperref}
\setlength{\parindent}{0pt}
\setlength{\parskip}{5pt}
\setlength{\emergencystretch}{3em}
\newcommand{\N}{\mathbb{N}}
\newcommand{\Z}{\mathbb{Z}}
\newcommand{\Q}{\mathbb{Q}}
\newcommand{\R}{\mathbb{R}}
\newcommand{\C}{\mathbb{C}}
\newcommand{\F}{\mathbb{F}}
\newcommand{\A}{\mathbb{A}}
\newcommand{\PP}{\mathbb{P}}
\newcommand{\E}{\mathbb{E}}
\newcommand{\eps}{\varepsilon}
\newcommand{\dd}{\,\mathrm{d}}
\begin{document}
\section*{UnsolvedMath ID: 701256}
\textbf{Source problem code:} catalogue entry
\subsection{Definitions and mathematical statement}
A \emph{group variety} is a class of groups defined by a collection of group identities. A variety \(\mathcal V\) is \emph{non-abelian} if it contains a non-abelian group. A non-abelian variety \(\mathcal V\) is called \emph{pseudo-abelian} if every finite group belonging to \(\mathcal V\) is abelian.

Write \(H\trianglelefteq G\) when \(H\) is a normal subgroup of \(G\). A group variety \(\mathcal V\) is a \emph{\(t\)-variety} if normality is transitive within every group in \(\mathcal V\); explicitly, for every \(G\in\mathcal V\) and all subgroups \(K\leq H\leq G\),
\[
K\trianglelefteq H\ \text{ and }\ H\trianglelefteq G
\quad\Longrightarrow\quad K\trianglelefteq G.
\]

\textbf{Question (O.~Macedo\'{n}ska, 2006).} Is every pseudo-abelian group variety a \(t\)-variety? Equivalently, if a group variety \(\mathcal V\) contains a non-abelian group and all its finite groups are abelian, must the displayed implication hold for every \(G\in\mathcal V\) and every subgroup chain \(K\leq H\leq G\)?

The supplied source reports that every non-abelian \(t\)-variety is pseudo-abelian, citing O.~Macedo\'{n}ska and A.~Storozhev (1997), and that the pseudo-abelian varieties constructed by A.~Yu.~Olshanskii (1986) are \(t\)-varieties. These background assertions are retained as source reports and have not been independently verified. The supplied registry labels the question open; this label is not an independent assessment of its current status.
\subsection{Short English statement}
A pseudo-abelian variety contains a non-abelian group, but every finite group in it is abelian. Must every group in such a variety have transitive normality: whenever \(K\) is normal in \(H\) and \(H\) is normal in \(G\), is \(K\) normal in \(G\)?
\subsection{Sources}
\begin{itemize}
\item Local UnsolvedMath registry, supplied source record: permanent ID \(701256\), exact source problem code \texttt{catalogue entry}, title \emph{Kourovka Notebook Problem 16.64}. The record attributes the problem to O.~Macedo\'{n}ska, proposed in 2006, and labels it open.
\item \emph{The Kourovka Notebook (Unsolved Problems in Group Theory)}, original problem 16.64, 16th issue (2006); the supplied record also identifies the 21st issue edition (2026), \texttt{arXiv:1401.0300v48}. The collection description names E.~I.~Khukhro and V.~D.~Mazurov as editors.
\item O.~Macedo\'{n}ska and A.~Storozhev, \emph{Commun. Algebra} \textbf{25}, no.~5 (1997), 1589--1593. Supplied reference for the assertion that every non-abelian \(t\)-variety is pseudo-abelian; no article title is supplied.
\item A.~Yu.~Olshanskii, \emph{Math. USSR--Sb.} \textbf{54} (1986), 57--80. Supplied reference for constructions of pseudo-abelian varieties which the source states are \(t\)-varieties; no article title is supplied.
\end{itemize}
\end{document}
